% GENERATED FILE. Edit canonical lessons, then run npm run generate:books.
% canonical-source-hash: fnv1a64:3e9458587dc499a8
% canonical-lessons: BN-C430-khet, BN-C430-langol, BN-C430-krishok, BN-C430-sech-deoya, BN-C430-chash-kora, BN-R430-review-424-430

\chapter{Fields and Farming}
\label{ch:bn-fields-and-farming}
\hlchaptermodality{430}

\begin{chapteropening}
\textbf{By the end of this chapter:} \emph{I can talk about a farmer, his plough and his fields, and watering the crop.}

Complete the last lesson of chapter 430: I can talk about a farmer, his plough and his fields, and watering the crop.
\end{chapteropening}

\section[\texorpdfstring{\bn{খেত} (khet)}{খেত (khet)}]{\bn{খেত} (khet) — a field (of crops)}

\label{lesson:BN-C430-khet}



\begin{quote}
Before the new one: say the Bengali for to shake (something), then the Bengali for to bend low.
\end{quote}

\subsection*{You'll want to know: \bn{খেত}}
\textbf{\bn{খেত}} — \emph{khet} — \textquotedblleft{}a field (of crops)\textquotedblright{}.

\begin{quote}
\bn{গ্রামের} \bn{চারপাশে} \bn{ধানের} \bn{খেত।}
\end{quote}

\begin{quote}
\emph{grāmer chārpāshe dhāner khet.}
\end{quote}

\begin{quote}
All round the village are rice fields.
\end{quote}

\bn{মাঠ} is open ground; \bn{খেত} is ground that grows something: \bn{ধানের} \bn{খেত}, a rice field.

\begin{cousinweb}
\textbf{\bn{খেত}} (\emph{khet}) is inherited: CDIAL lists Bengali \emph{khet}, \textquotedblleft{}a field\textquotedblright{}, under Sanskrit \emph{kṣ\'{\={e}}tra}, \textquotedblleft{}land\textquotedblright{} (CDIAL 3735).
\end{cousinweb}

\begin{grammarlens}[title={What you've built}]
One more word for getting around.
\end{grammarlens}

\subsection*{Your turn}
\begin{itemize}
\raggedright
  \item \emph{Say it:} \emph{khet}
  \item \emph{Say it:} \emph{khet}, once more
  \item \emph{Say it:} \emph{grāmer chārpāshe dhāner khet.} — All round the village are rice fields.
  \item \emph{Recall:} say the Bengali for to shake (something), then the Bengali for to bend low, then say \emph{khet} again
\end{itemize}

\begin{tcolorbox}[breakable,colback=teal!4,colframe=teal!35!black,title={Before you move on}]
What does \bn{খেত} mean? (\textquotedblleft{}A field (of crops)\textquotedblright{}.) Say it once more.
\end{tcolorbox}
\section[\texorpdfstring{\bn{লাঙল} (lāngôl)}{লাঙল (lāngôl)}]{\bn{লাঙল} (lāngôl) — a plough}

\label{lesson:BN-C430-langol}



\begin{quote}
Before the new one: say the Bengali for to bend low, then the Bengali for a field (of crops).
\end{quote}

\subsection*{You'll want to know: \bn{লাঙল}}
\textbf{\bn{লাঙল}} — \emph{lāngôl} — \textquotedblleft{}a plough\textquotedblright{}.

\begin{quote}
\bn{গরু} \bn{লাঙল} \bn{টানে।}
\end{quote}

\begin{quote}
\emph{gôru lāngôl ṭāne.}
\end{quote}

\begin{quote}
The cattle pull the plough.
\end{quote}

\begin{cousinweb}
\textbf{\bn{লাঙল}} (\emph{lāngôl}) is inherited: CDIAL lists Bengali \emph{lāṅal}, \textquotedblleft{}a plough\textquotedblright{}, under Sanskrit \emph{l\'{\={a}}ṅgala}, \textquotedblleft{}a plough\textquotedblright{} (CDIAL 11006).
\end{cousinweb}

\begin{grammarlens}[title={What you've built}]
One more word for getting around.
\end{grammarlens}

\subsection*{Your turn}
\begin{itemize}
\raggedright
  \item \emph{Say it:} \emph{lāngôl}
  \item \emph{Say it:} \emph{lāngôl}, once more
  \item \emph{Say it:} \emph{gôru lāngôl ṭāne.} — The cattle pull the plough.
  \item \emph{Recall:} say the Bengali for to bend low, then the Bengali for a field (of crops), then say \emph{lāngôl} again
\end{itemize}

\begin{tcolorbox}[breakable,colback=teal!4,colframe=teal!35!black,title={Before you move on}]
What does \bn{লাঙল} mean? (\textquotedblleft{}A plough\textquotedblright{}.) Say it once more.
\end{tcolorbox}
\section[\texorpdfstring{\bn{কৃষক} (krishôk)}{কৃষক (krishôk)}]{\bn{কৃষক} (krishôk) — a farmer}

\label{lesson:BN-C430-krishok}



\begin{quote}
Before the new one: say the Bengali for a field (of crops), then the Bengali for a plough.
\end{quote}

\subsection*{You'll want to know: \bn{কৃষক}}
\textbf{\bn{কৃষক}} — \emph{krishôk} — \textquotedblleft{}a farmer\textquotedblright{}.

\begin{quote}
\bn{কৃষক} \bn{ভোরে} \bn{খেতে} \bn{যান।}
\end{quote}

\begin{quote}
\emph{krishôk bhore khete jān.}
\end{quote}

\begin{quote}
The farmer goes to the field at dawn.
\end{quote}

A careful word, common in the news and in schoolbooks.

\begin{cousinweb}
\textbf{\bn{কৃষক}} (\emph{krishôk}) is Sanskrit \emph{kṛṣaka}, \textquotedblleft{}a ploughman, a husbandman\textquotedblright{} (Monier-Williams), from \emph{kṛṣ}, to plough.
\end{cousinweb}

\begin{grammarlens}[title={What you've built}]
One more word for getting around.
\end{grammarlens}

\subsection*{Your turn}
\begin{itemize}
\raggedright
  \item \emph{Say it:} \emph{krishôk}
  \item \emph{Say it:} \emph{krishôk}, once more
  \item \emph{Say it:} \emph{krishôk bhore khete jān.} — The farmer goes to the field at dawn.
  \item \emph{Recall:} say the Bengali for a field (of crops), then the Bengali for a plough, then say \emph{krishôk} again
\end{itemize}

\begin{tcolorbox}[breakable,colback=teal!4,colframe=teal!35!black,title={Before you move on}]
What does \bn{কৃষক} mean? (\textquotedblleft{}A farmer\textquotedblright{}.) Say it once more.
\end{tcolorbox}
\section[\texorpdfstring{\bn{সেচ} \bn{দেওয়া} (sech deoyā)}{সেচ দেওয়া (sech deoyā)}]{\bn{সেচ} \bn{দেওয়া} (sech deoyā) — to irrigate; to water (fields)}

\label{lesson:BN-C430-sech-deoya}



\begin{quote}
Before the new one: say the Bengali for a plough, then the Bengali for a farmer.
\end{quote}

\subsection*{You'll want to know: \bn{সেচ} \bn{দেওয়া}}
\textbf{\bn{সেচ} \bn{দেওয়া}} — \emph{sech deoyā} — \textquotedblleft{}to irrigate; to water (fields)\textquotedblright{}.

\begin{quote}
\bn{শুকনো} \bn{দিনে} \bn{কৃষক} \bn{খেতে} \bn{সেচ} \bn{দেন।}
\end{quote}

\begin{quote}
\emph{shukno dine krishôk khete sech den.}
\end{quote}

\begin{quote}
On dry days the farmer irrigates the field.
\end{quote}

Past (he/she …): \bn{তিনি} \bn{খেতে} \bn{সেচ} \bn{দিলেন}, \emph{tini khete sech dilen}.

\begin{cousinweb}
\textbf{\bn{সেচ} \bn{দেওয়া}} (\emph{sech deoyā}) joins \bn{সেচ}, \textquotedblleft{}irrigation\textquotedblright{}, of the family CDIAL sets under Sanskrit \emph{siñcáti}, \textquotedblleft{}pours out, sprinkles\textquotedblright{}, with Bengali \emph{sẽcā}, \textquotedblleft{}to irrigate, to bale out\textquotedblright{} (CDIAL 13394), to \bn{দেওয়া}, to give.
\end{cousinweb}

\begin{grammarlens}[title={What you've built}]
One more word for getting around.
\end{grammarlens}

\subsection*{Your turn}
\begin{itemize}
\raggedright
  \item \emph{Say it:} \emph{sech deoyā}
  \item \emph{Say it:} \emph{sech deoyā}, once more
  \item \emph{Say it:} \emph{shukno dine krishôk khete sech den.} — On dry days the farmer irrigates the field.
  \item \emph{Recall:} say the Bengali for a plough, then the Bengali for a farmer, then say \emph{sech deoyā} again
\end{itemize}

\begin{tcolorbox}[breakable,colback=teal!4,colframe=teal!35!black,title={Before you move on}]
What does \bn{সেচ} \bn{দেওয়া} mean? (\textquotedblleft{}To irrigate; to water (fields)\textquotedblright{}.) Say it once more.
\end{tcolorbox}
\section[\texorpdfstring{\bn{চাষ} \bn{করা} (chāsh kôrā)}{চাষ করা (chāsh kôrā)}]{\bn{চাষ} \bn{করা} (chāsh kôrā) — to farm; to grow (a crop)}

\label{lesson:BN-C430-chash-kora}



\begin{quote}
Before the new one: say the Bengali for a farmer, then the Bengali for to irrigate.
\end{quote}

\subsection*{You'll want to know: \bn{চাষ} \bn{করা}}
\textbf{\bn{চাষ} \bn{করা}} — \emph{chāsh kôrā} — \textquotedblleft{}to farm; to grow (a crop)\textquotedblright{}.

\begin{quote}
\bn{এই} \bn{গ্রামের} \bn{লোকেরা} \bn{ধান} \bn{চাষ} \bn{করে।}
\end{quote}

\begin{quote}
\emph{ei grāmer lokerā dhān chāsh kôre.}
\end{quote}

\begin{quote}
The people of this village grow rice.
\end{quote}

Past (he/she …): \bn{তিনি} \bn{জমিতে} \bn{সবজি} \bn{চাষ} \bn{করলেন}, \emph{tini jômite sôbji chāsh kôrlen}.

\bn{চাষ} alone is farming, the work of the fields.

\begin{cousinweb}
\textbf{\bn{চাষ} \bn{করা}} (\emph{chāsh kôrā}) joins \bn{চাষ}, which CDIAL lists as Bengali \emph{cās}, \textquotedblleft{}ploughing\textquotedblright{}, under the reconstructed \emph{carṣati}, \textquotedblleft{}ploughs\textquotedblright{} (CDIAL 4712), to \bn{করা}, to do.
\end{cousinweb}

\begin{grammarlens}[title={What you've built}]
One more word for getting around.
\end{grammarlens}

\subsection*{Your turn}
\begin{itemize}
\raggedright
  \item \emph{Say it:} \emph{chāsh kôrā}
  \item \emph{Say it:} \emph{chāsh kôrā}, once more
  \item \emph{Say it:} \emph{ei grāmer lokerā dhān chāsh kôre.} — The people of this village grow rice.
  \item \emph{Recall:} say the Bengali for a farmer, then the Bengali for to irrigate, then say \emph{chāsh kôrā} again
\end{itemize}

\begin{tcolorbox}[breakable,colback=teal!4,colframe=teal!35!black,title={Before you move on}]
What does \bn{চাষ} \bn{করা} mean? (\textquotedblleft{}To farm; to grow (a crop)\textquotedblright{}.) Say it once more.
\end{tcolorbox}
\section[(dialogue)]{Looking back: from Prices Up and Down to Fields and Farming}

\label{lesson:BN-R430-review-424-430}



\begin{quote}
Say the Bengali for a sale (at reduced prices) and an offer.
\end{quote}

\subsection*{Your turn}
Say these aloud:

\begin{itemize}
\raggedright
  \item a sale (at reduced prices), an offer, a slump, demand (for goods), to fluctuate
  \item a cloud, the monsoon, in torrents, a lightning strike, to spread over
  \item cold (to the touch), snow, to freeze, bitterly cold, to melt
  \item summer, autumn, a season (for a crop or an activity), temperature, heat (warmth)
  \item humid, muggy, hazy, damp (and dank), to be covered over
  \item a nor'wester, hail, to roar, to shake (something), to bend low
  \item a field (of crops), a plough, a farmer, to irrigate, to farm
\end{itemize}

\begin{tcolorbox}[breakable,colback=teal!4,colframe=teal!35!black,title={Before you move on}]
A sale (at reduced prices)? (\textbf{\bn{সেল}}, \emph{sel}.) An offer? (\textbf{\bn{অফার}}, \emph{ôphār}.) A slump? (\textbf{\bn{মন্দা}}, \emph{môndā}.) Demand (for goods)? (\textbf{\bn{চাহিদা}}, \emph{chāhidā}.) To fluctuate? (\textbf{\bn{ওঠানামা} \bn{করা}}, \emph{oṭhānāmā kôrā}.) A cloud? (\textbf{\bn{মেঘ}}, \emph{megh}.) The monsoon? (\textbf{\bn{বর্ষা}}, \emph{bôrshā}.) In torrents? (\textbf{\bn{মুষলধারে}}, \emph{mushôldhāre}.) A lightning strike? (\textbf{\bn{বজ্রপাত}}, \emph{bôjrôpāt}.) To spread over? (\textbf{\bn{ছাওয়া}}, \emph{chhāoyā}.) Cold (to the touch)? (\textbf{\bn{ঠান্ডা}}, \emph{ṭhānḍā}.) Snow? (\textbf{\bn{তুষার}}, \emph{tushār}.) To freeze? (\textbf{\bn{জমে} \bn{যাওয়া}}, \emph{jôme jāoyā}.) Bitterly cold? (\textbf{\bn{কনকনে}}, \emph{kônkône}.) To melt? (\textbf{\bn{গলে} \bn{যাওয়া}}, \emph{gôle jāoyā}.) Summer? (\textbf{\bn{গ্রীষ্ম}}, \emph{grīshmô}.) Autumn? (\textbf{\bn{শরৎ}}, \emph{shôrôt}.) A season (for a crop or an activity)? (\textbf{\bn{মৌসুম}}, \emph{mousum}.) Temperature? (\textbf{\bn{তাপমাত্রা}}, \emph{tāpmātrā}.) Heat (warmth)? (\textbf{\bn{উত্তাপ}}, \emph{uttāp}.) Humid? (\textbf{\bn{আর্দ্র}}, \emph{ārdrô}.) Muggy? (\textbf{\bn{ভ্যাপসা}}, \emph{bhæpsā}.) Hazy? (\textbf{\bn{ঝাপসা}}, \emph{jhāpsā}.) Damp (and dank)? (\textbf{\bn{স্যাঁতসেঁতে}}, \emph{sæ̃tsẽte}.) To be covered over? (\textbf{\bn{ঢেকে} \bn{যাওয়া}}, \emph{ḍheke jāoyā}.) A nor'wester? (\textbf{\bn{কালবৈশাখী}}, \emph{kālboishākhī}.) Hail? (\textbf{\bn{শিলাবৃষ্টি}}, \emph{shilābrishṭi}.) To roar? (\textbf{\bn{গর্জানো}}, \emph{gôrjāno}.) To shake (something)? (\textbf{\bn{কাঁপানো}}, \emph{kā̃pāno}.) To bend low? (\textbf{\bn{নুয়ে} \bn{পড়া}}, \emph{nuye pôṛā}.) A field (of crops)? (\textbf{\bn{খেত}}, \emph{khet}.) A plough? (\textbf{\bn{লাঙল}}, \emph{lāngôl}.) A farmer? (\textbf{\bn{কৃষক}}, \emph{krishôk}.) To irrigate? (\textbf{\bn{সেচ} \bn{দেওয়া}}, \emph{sech deoyā}.) To farm? (\textbf{\bn{চাষ} \bn{করা}}, \emph{chāsh kôrā}.)
\end{tcolorbox}
